\subsection{RQ1.5: AKS Transfer Validation}
\label{sec:analysis-rq15}

The results in \cref{sec:rq15} reveal four bounded findings for the AKS
transfer-validation stage.

\paragraph{Finding 1: AKS preserves the local Kubernetes condition ordering.}
The complete RQ1.4 ordering is preserved in AKS: the monolithic Kubernetes
service is fastest, followed by \texttt{chain\_2svc}, \texttt{chain\_3svc},
\texttt{chain\_4svc}, and \texttt{chain\_5svc}. This provides directional
transfer evidence. The same service-count trend that increased latency in
local Kubernetes remains visible under managed cloud execution, even though
the absolute millisecond values differ between environments.

\paragraph{Finding 2: AKS overhead is monotonic with service count.}
Relative to the AKS monolithic baseline, mean latency rises from 70.802~ms to
73.437~ms, 75.410~ms, 78.294~ms, and 80.519~ms as the chain grows from one to
five services. The corresponding overheads are +2.635~ms (3.7\%), +4.608~ms
(6.5\%), +7.492~ms (10.6\%), and +9.717~ms (13.7\%). Within this AKS
deployment, adding synchronous service boundaries is therefore associated with
measurable end-to-end cost.

\paragraph{Finding 3: Marginal overhead is bounded by the activation profile.}
The AKS marginal increments are +2.635~ms from the monolithic baseline to
\texttt{chain\_2svc}, +1.973~ms from \texttt{chain\_2svc} to
\texttt{chain\_3svc}, +2.884~ms from \texttt{chain\_3svc} to
\texttt{chain\_4svc}, and +2.225~ms from \texttt{chain\_4svc} to
\texttt{chain\_5svc}. The first three steps reflect the cost of adding
synchronous boundaries that move 392--784~KiB activations, while the final
step is smaller than the preceding one because the layer-4 boundary moves
only 98~KiB. This is consistent with the bounded non-linearity expected from
the activation profile of the chain family: an additional service boundary
does not imply a constant per-boundary penalty.

\paragraph{Finding 4: Absolute latency is environment-specific.}
All AKS means are lower than their local Kubernetes counterparts, including
the monolithic baseline. The lower AKS absolute latencies should not be
interpreted as evidence that AKS is generally faster than local Kubernetes.
The local and AKS runs differ in host CPU, kernel, virtualisation layer,
Kubernetes implementation, container runtime context, CPU scheduling context,
and networking path. RQ1.5 therefore uses within-environment ordering and
normalised overhead as the primary comparison, not direct absolute latency. On
that basis, AKS provides directional transfer evidence for the RQ1.4 ordering
while preserving the claim scope: absolute latency is deployment-specific, and
the AKS run should not be interpreted as a numerical replication of the local
Kubernetes measurements.

\paragraph{Answer to RQ1.5.}
Under the controlled AKS transfer-validation setup, the local Kubernetes
service-count pattern transfers directionally to managed cloud execution. The
AKS run preserves the complete RQ1.4 ordering: \texttt{monolithic\_k8s\_1svc}
remains fastest at 70.802~ms, followed by \texttt{chain\_2svc} at 73.437~ms,
\texttt{chain\_3svc} at 75.410~ms, \texttt{chain\_4svc} at 78.294~ms, and
\texttt{chain\_5svc} at 80.519~ms. Relative overhead grows monotonically from
3.7\% for \texttt{chain\_2svc} to 13.7\% for \texttt{chain\_5svc}, with the
final marginal increment from \texttt{chain\_4svc} to \texttt{chain\_5svc}
(+2.225~ms) smaller than the preceding step (+2.884~ms), consistent with the
activation-transfer-size explanation from RQ1.3. RQ1.5 therefore supports a
directional transfer claim: the same architectural chain ordering remains
visible in AKS, although absolute latency values are environment-specific and
should not be interpreted as a direct replication of the local Kubernetes
measurements.
